% TOP_PROBLEM: {"problemId": "problem.witsenhausen-counterexample-optimal-strategy-and-cost", "canonicalTitle": "Witsenhausen Counterexample: Optimal Strategy and Cost", "releaseRank": 488, "primaryDomain": "applied_computational_mathematics", "primaryDomainLabel": "Applied & computational mathematics", "displayStatus": "open", "releaseStatus": "open"}
% !TeX program = lualatex
% ATTEMPT_STATUS: unresolved
\documentclass[11pt]{article}
\usepackage[margin=25mm]{geometry}
\usepackage{fontspec}
\setmainfont{Latin Modern Roman}
\usepackage{amsmath,amssymb,mathtools}
\usepackage{unicode-math}
\setmathfont{Latin Modern Math}
\usepackage{xurl,hyperref}
\hypersetup{colorlinks=true,urlcolor=blue}
\setcounter{secnumdepth}{0}
\setlength{\parindent}{0pt}
\setlength{\parskip}{5pt}
\setlength{\emergencystretch}{3em}
\begin{document}
\section*{\texorpdfstring{Witsenhausen Counterexample: Optimal Strategy and Cost}{Witsenhausen Counterexample: Optimal Strategy and Cost}}\label{488}
\subsection{Definitions and mathematical statement}
Fix the usual parameters $k>0$ and $\sigma>0$. Let $X_0\sim N(0,\sigma^2)$ and $V\sim N(0,1)$ be independent. The first decision maker chooses $U_1=\gamma_1(X_0)$, after which
\[
 X_1=X_0+U_1,\qquad Z_2=X_1+V.
\]
The second chooses $U_2=\gamma_2(Z_2)$ and $X_2=X_1+U_2$. Over measurable strategies with finite expected cost, determine the exact infimum and optimal strategies for
\[
 J^*(k,\sigma)=\inf_{\gamma_1,\gamma_2}
 {\mathbb{E}}\left[k\gamma_1(X_0)^2+
       (X_0+\gamma_1(X_0)+\gamma_2(Z_2))^2\right].
\]
The first control changes what the second decision maker observes; this information structure is the essential constraint. Arbitrary nonlinear strategies are allowed. The catalogue uses $kU_1^2$, not the common alternative notation $k^2U_1^2$; changing that convention requires reparametrizing $k$.
\par\smallskip\textit{Formulation and concept references:} \href{https://people.seas.harvard.edu/~ho/DEDS/OO/Lecture4/Slide01.html}{[S1]}.

\subsection{Short English statement}
What are the exact optimal nonlinear policies and minimum expected cost in Witsenhausen's two-stage control problem with its specified information structure?
\subsection{Research attempt}
\paragraph{Scope and source check.}
This attempt by GPT-6 Astra Ultra was prepared on 27 September 2026.
Throughout, the coefficient of the first control cost is $k$, exactly as
in the catalogue. A source using $k_0^2U_1^2$ corresponds to $k=k_0^2$.
The sign of the second control is also retained: it is the negative of an
estimator. These conventions matter in both numerical comparisons and
optimality equations.

The historical slide [S1] illustrates nonlinear signalling. Park, Grover,
and Sahai [E1] give a uniform constant-factor approximation, which is
different from the exact optimum requested here. The recent manuscript
[E2] uses a transport formulation and studies stationarity and numerical
quantization. The reduction below is proved directly; none of its
additional regularity or asymptotic claims is needed. In particular, this
attempt does not identify a minimizing law for every pair of parameters.
The concrete progress consists of an exact one-controller reduction,
a rigorous separation from every affine policy at one parameter pair,
and a universal scalar lower bound.

\paragraph{Lemma 488.1: eliminate the second controller exactly.}
Write $X=X_0$, $g(x)=x+\gamma_1(x)$, $S=g(X)$, and $Y=S+V$.
Finite first-stage cost implies $S\in L^2$. Let $m(Y)=\mathbb E[S\mid Y]$.
For every admissible second controller $h$, orthogonal projection gives
\begin{equation}\label{a488:projection}
 \mathbb E(S+h(Y))^2
 =\mathbb E(S-m(Y))^2+\mathbb E(m(Y)+h(Y))^2.
\end{equation}
Indeed, the cross term has conditional expectation zero given $Y$.
If the left side is finite, $h(Y)\in L^2$ by the triangle inequality;
therefore the calculation covers all finite-cost competitors. The unique
best second action, up to observation-null sets, is $h=-m$.

For $\phi(t)=(2\pi)^{-1/2}e^{-t^2/2}$ and
$\phi_\sigma(x)=\sigma^{-1}\phi(x/\sigma)$, a version is
\begin{equation}\label{a488:posterior}
 m_g(y)=
 \frac{\int_{\mathbb R}g(x)\phi(y-g(x))\phi_\sigma(x)\,dx}
      {\int_{\mathbb R}\phi(y-g(x))\phi_\sigma(x)\,dx}.
\end{equation}
The denominator is strictly positive. The numerator is finite since
$\phi$ is bounded and $\mathbb E|S|<\infty$. Thus the exact target is
\begin{equation}\label{a488:reduced}
 J^*(k,\sigma)=\inf_{g:\,\mathbb E(g(X)-X)^2<\infty}
 \left\{k\mathbb E(g(X)-X)^2+
              \mathbb E(S-m_g(S+V))^2\right\}.
\end{equation}
This is a variational reformulation, not an evaluation of the infimum.
It removes one information constraint while leaving a nonlinear
functional of $g$ through the full observation distribution.

\paragraph{A useful normalization that does not assume oddness.}
Replacing $S$ by $S+c$ leaves its estimation error unchanged: translate
the observation and then translate the posterior mean. Meanwhile
\[
 \mathbb E(S+c-X)^2
 =\mathbb E(S-X)^2+2c\mathbb E(S-X)+c^2.
\]
Choosing $c=-\mathbb E(S-X)$ cannot increase the objective. Hence the
infimum may be restricted to $\mathbb ES=0$. This argument does not show
that $g$ is odd, or that an optimizer is a uniform quantizer. Symmetry of
the Gaussian input makes reflection a symmetry of the problem, but
averaging two reflected policies changes their signalling distributions.
Convexity needed to justify that averaging is not available here.

\paragraph{Lemma 488.2: an exact optimization over output laws.}
Let $\mu=N(0,\sigma^2)$ and let $\nu$ range over probability laws on
$\mathbb R$ with finite second moment. Define
\[
 W_2^2(\mu,\nu)=\inf\mathbb E(A-B)^2,
 \qquad A\sim\mu,\ B\sim\nu,
\]
where the infimum ranges over couplings. Let
$D(\nu)=\mathbb E(S-\mathbb E[S\mid S+V])^2$ for $S\sim\nu$
independent of $V$. Then
\begin{equation}\label{a488:transport}
 J^*(k,\sigma)=\inf_{\nu}\{kW_2^2(\mu,\nu)+D(\nu)\}.
\end{equation}
For a fixed law $\nu$, the estimation term depends only on that law.
The least transport cost is attained by the increasing quantile map
\[
 T_\nu(x)=F_\nu^{-1}(F_\mu(x)),\qquad
 F_\nu^{-1}(u)=\inf\{t:F_\nu(t)\ge u\}.
\]
Here $F_\mu(X)$ is uniform on $(0,1)$, so this is a deterministic
admissible controller even when $\nu$ has atoms.

For completeness, the optimal-coupling fact follows by maximizing
$\mathbb E(AB)$ with fixed marginals. For bounded nonnegative variables,
the layer-cake identity expresses that expectation as
\[
 \int_0^\infty\!\int_0^\infty
        \mathbb P(A>s,B>t)\,ds\,dt.
\]
Each integrand is at most the smaller marginal tail probability. The
common-uniform quantile coupling attains all these bounds simultaneously
because its upper-tail events are nested intervals in $(0,1)$. Translating
bounded variables makes them nonnegative without changing which coupling
maximizes the product. Truncation and $L^2$ convergence extend the result
to arbitrary finite second moments. Expanding $(A-B)^2$ proves the claim.

Consequently one may restrict to nondecreasing $g$ without increasing the
infimum. This is more informative than merely centering $g$, but it does
not determine $\nu$: smooth, atomic, and mixed laws are still competitors.
We make no assertion here that a minimizing law must have any one of
these forms. Equation \eqref{a488:transport} agrees with the type of
reformulation in [E2], with the parameter conversion already specified.

\paragraph{Lemma 488.3: the affine benchmark is a one-variable problem.}
Take $g(x)=\alpha x+b$. Gaussian conditioning gives
\[
 m_g(y)=b+\frac{\alpha^2\sigma^2}{1+\alpha^2\sigma^2}(y-b),
 \qquad
 D=\frac{\alpha^2\sigma^2}{1+\alpha^2\sigma^2}.
\]
Hence the best cost with affine first control, even allowing an arbitrary
second controller, is
\begin{equation}\label{a488:affine}
 J_{\rm aff}(k,\sigma)
 =\min_{0\le\alpha\le1}
 \left\{k\sigma^2(\alpha-1)^2+
                \frac{\alpha^2\sigma^2}{1+\alpha^2\sigma^2}\right\}.
\end{equation}
The optimum has $b=0$. Replacing $\alpha<0$ by $|\alpha|$ lowers the
first term and preserves the second; replacing $\alpha>1$ by $1$ lowers
both. Endpoints have derivative of opposite nonzero signs, so an optimum
lies in $(0,1)$ and satisfies
\begin{equation}\label{a488:quintic}
       k(\alpha-1)(1+\sigma^2\alpha^2)^2+\alpha=0.
\end{equation}
There may be several stationary roots. Solving one root or initializing
near $\alpha=1$ is insufficient; every admissible root must be compared.
The trivial policies also give the universal upper bound
\[
      J^*(k,\sigma)\le\min\{k\sigma^2,\sigma^2/(1+\sigma^2)\}.
\]

\paragraph{Proposition 488.4: a certified nonlinear improvement.}
For $a>0$, choose $g(x)=a\operatorname{sgn}(x)$ and
$h(y)=-a\operatorname{sgn}(y)$; values at zero do not matter.
Let $Q(a)=\int_a^\infty\phi(t)\,dt$. Direct calculation gives
\begin{equation}\label{a488:signcost}
 J_{\rm sign}(a)=k\left(\sigma^2+a^2-2a\sigma\sqrt{2/\pi}\right)
                        +4a^2Q(a).
\end{equation}
The first identity uses $\mathbb E|X|=\sigma\sqrt{2/\pi}$.
Conditional on $S=a$, the decoder is wrong exactly when $V<-a$;
the other sign is symmetric. An error has squared residual $4a^2$.
This also proves finiteness of the policy cost.

Set $k=1/25$, $\sigma=5$, and $a=4$. Every affine first policy has
cost at least $9/16$. In fact, for $\alpha\le1/4$ its first term
$(\alpha-1)^2$ is at least $9/16$, while for $\alpha\ge1/4$ its
second term is at least $25/41>9/16$. An intercept only adds cost.
On the other hand,
\[
 J_{\rm sign}(4)=\frac{41}{25}-\frac85\sqrt{2/\pi}+64Q(4)
                                      <\frac38.
\]
Here is a rationally checkable bound, independent of numerical quadrature.
Since $\pi<22/7$, one has $\sqrt{2/\pi}>797/1000$.
For $a>0$,
\[
 Q(a)\le\frac1a\int_a^\infty t\phi(t)\,dt
                  =\frac{e^{-a^2/2}}{a\sqrt{2\pi}}.
\]
The finite sum $\sum_{j=0}^{14}8^j/j!>2900$ and $\sqrt{2\pi}>2$
give $64Q(4)<8/2900$. Finally
\[
 \frac{41}{25}-\frac85\frac{797}{1000}+\frac8{2900}<\frac38.
\]
Thus $J^*(1/25,5)<3/8<9/16\le J_{\rm aff}(1/25,5)$.
This rederives the basic nonlinear advantage with deliberately loose,
transparent constants. It does not claim that the two-level policy is
globally optimal or that these are the best known bounds.

The best decoder for this same two-level encoder improves the bound:
Bayes' rule gives $m(y)=a\tanh(ay)$ and posterior variance
$a^2\operatorname{sech}^2(ay)$. The hard sign decoder is useful because
its performance is especially easy to certify; optimizing the decoder
does not certify the encoder.

\paragraph{Lemma 488.5: a lower bound that holds for every measurable policy.}
Let $P=\mathbb E(S-X)^2$ and let $D=D(\mathcal L(S))$. Then
\begin{equation}\label{a488:lower}
 J^*(k,\sigma)\ge\inf_{P\ge0}
 \left\{kP+
  \left(\frac{\sigma}{\sqrt{1+(\sigma+\sqrt P)^2}}-\sqrt P\right)_+^2
 \right\},\qquad r_+=\max\{r,0\}.
\end{equation}
To prove this, $\mathbb ES^2\le(\sigma+\sqrt P)^2$ by the $L^2$
triangle inequality. For the additive Gaussian channel, with information
measured in natural logarithms,
\[
 I(X;Y)=I(S;Y)\le\tfrac12\log(1+\operatorname{Var}S)
                         \le\tfrac12\log(1+(\sigma+\sqrt P)^2).
\]
The equality follows since $S$ is a function of $X$ and the conditional
law of $Y$ given $X$ depends only on $S$. The upper bound follows by
maximizing entropy at fixed variance: $h(Y)\le\frac12\log(2\pi e
\operatorname{Var}Y)$ and $h(Y\mid X)=\frac12\log(2\pi e)$.

For any square-integrable estimator $\widehat X(Y)$ of the Gaussian
variable $X$, conditional entropy and the same variance bound imply
\[
 I(X;Y)\ge\tfrac12\log
            \frac{\sigma^2}{\mathbb E(X-\widehat X(Y))^2}.
\]
Indeed $h(X\mid Y)=h(X-\widehat X(Y)\mid Y)$ is at most the entropy
of the error, and that entropy is bounded by its second moment.
The entropy inequalities can equivalently be expressed using nonnegative
relative entropy; degenerate cases follow by regularization.
Applying the result to $\widehat X=m(Y)$ and then using
\[
 \|X-m(Y)\|_2\le\|X-S\|_2+\|S-m(Y)\|_2=\sqrt P+\sqrt D
\]
proves \eqref{a488:lower}. The expression on the right is strictly
positive for fixed $k,\sigma>0$: it is continuous, tends to infinity
as $P\to\infty$, and can vanish only if $P=0$, where its second term
is $\sigma^2/(1+\sigma^2)>0$.

This bound is not asserted to improve the literature. Its purpose is to
show precisely where a universal estimate loses information. It replaces
the actual output law by a variance bound and treats the two errors in
a triangle inequality. Their equality conditions need not coexist with
an optimal signalling policy. A positive lower bound and a good quantizer
upper bound leave an interval; they do not determine its exact value.

\paragraph{Attempt to close the gap by alternating best responses.}
For a fixed decoder $h$, the first-controller objective is pointwise in
$x$:
\[
       H_x(t)=k(t-x)^2+\int_{\mathbb R}(t+h(t+v))^2\phi(v)\,dv.
\]
If $h$ is bounded and continuous, this is continuous and coercive in $t$,
and a measurable choice of minimizer gives a best first response. The
measurability follows, for example, by taking the smallest minimizer:
compact sublevel sets and continuity allow its level sets to be tested
with infima over rational intervals. The selected response has finite
expected cost, since $H_x(t_{\min})\leq H_x(x)\leq
(|x|+\|h\|_\infty)^2$, which is integrable under the Gaussian prior.
Under differentiability and
domination sufficient to differentiate under the integral, an interior
stationary response satisfies
\begin{equation}\label{a488:firstresponse}
 k(t-x)+\mathbb E[(t+h(t+V))(1+h'(t+V))]=0.
\end{equation}
The bounded-continuous argument covers the posterior decoder of the
two-level policy. It is not an existence proof for best responses to
every arbitrary measurable decoder.

Alternating \eqref{a488:posterior} with minimization of $H_x$ makes a
useful search method. A fixed point, even one using global minimization
in each coordinate separately, need not minimize jointly. The elementary
function
\[
 F(s,t)=(st-1)^2+\varepsilon(s^2+t^2),\qquad0<\varepsilon<1/2,
\]
has $(0,0)$ as a pair of global coordinate best responses, with value
$1$, while $F(1,1)=2\varepsilon<1$. This is a diagnostic of the
logical implication, not a counterexample policy for this control problem.
No convexity theorem has been proved here that excludes such a failure
for the reduced Witsenhausen functional.

\paragraph{Verification and exact remaining obligation.}
The displayed nonlinear separation was checked with rational arithmetic
for the finite exponential sum and the two final comparisons; floating
point values are unnecessary to its proof. The posterior, affine, and
transport reductions are symbolic arguments, not results of a finite
search. The information lower bound is uniform over admissible policies.
None of these checks explores all output laws or establishes global
optimality of a numerical stationary point.

\textbf{UNRESOLVED.} To finish the catalogue problem one must identify,
for arbitrary $k,\sigma>0$, a law or policy attaining a value that a
matching lower bound proves minimal. The exact reductions and the proved
nonlinear-versus-affine separation do not supply that matching certificate.
No general solution or novelty claim is made.

\subsection{Sources}
\begin{itemize}
\item[S1]Formal statement, Lecture 4 slide, 'Witsenhausen Problem (1968)'.
\url{https://people.seas.harvard.edu/~ho/DEDS/OO/Lecture4/Slide01.html}.
\textit{Formulation source.}
\item[E1]Se Yong Park, Pulkit Grover, and Anant Sahai,
\textit{A constant-factor approximately optimal solution to the
Witsenhausen Counterexample}, author-hosted manuscript, Theorem 1.
\url{https://people.eecs.berkeley.edu/~sahai/Papers/CDC09Submission.pdf}.
Uses the coefficient $k^2$; cited for approximation context, not an exact optimum.
\item[E2]Quanyan Zhu, \textit{From Optimal Transport to Optimal
Quantization: A Variational Study of the Witsenhausen Counterexample},
arXiv:2609.18992v1, accessed 27 September 2026.
\url{https://arxiv.org/abs/2609.18992v1}.
The transport formulation is rederived here; additional manuscript claims
are not used as established premises.
\end{itemize}
\end{document}
